% GENERATED FILE. Edit canonical lessons, then run npm run generate:books.
% canonical-source-hash: fnv1a64:8dbcef828c0384d1
% canonical-lessons: MR-C301-disne, MR-C301-sapadne, MR-C301-adhalne, MR-C301-dokavne, MR-C301-nirakhne

\chapter{Strange Sights}
\label{ch:mr-strange-sights}
\hlchaptermodality{301}

\begin{chapteropening}
\textbf{By the end of this chapter:} \emph{I can say what someone saw, found or came across.}

Complete the last lesson of chapter 301: I can say what someone saw, found or came across.
\end{chapteropening}

\section[\texorpdfstring{\mr{दिसणे} (disṇe)}{दिसणे (disṇe)}]{\mr{दिसणे} (disṇe) — to be seen; to appear}

\label{lesson:MR-C301-disne}



\begin{quote}
Before the new one: say the Marathi for to flee, then the Marathi for to shut in.
\end{quote}

\subsection*{You'll want to know: \mr{दिसणे}}
\textbf{\mr{दिसणे}} — \emph{disṇe} — \textquotedblleft{}to be seen; to appear\textquotedblright{}.

\begin{quote}
\mr{खिडकीतून} \mr{मला} \mr{एक} \mr{विचित्र} \mr{प्रकाश} \mr{दिसला}.
\end{quote}

\begin{quote}
\emph{khiḍkītūn malā ek vicitra prakāś dislā.}
\end{quote}

\begin{quote}
I saw a strange light through the window.
\end{quote}

Past (to him …): \mr{त्याला} \mr{दिसलं}, \emph{tyālā dislaṁ}.

A dative-subject verb: \mr{मला} \mr{दिसलं}, I saw; the thing seen agrees: \mr{मला} \mr{घर} \mr{दिसलं}, \mr{मला} \mr{गाडी} \mr{दिसली}.

\begin{cousinweb}
\textbf{\mr{दिसणे}} (\emph{disṇe}) is a tadbhava: CDIAL lists Marathi \emph{disṇe}, 'to appear, seem', under Sanskrit \emph{dṛśyate}, 'is seen, appears' (CDIAL 6516).
\end{cousinweb}

\begin{grammarlens}[title={What you've built}]
One more word for telling what happened.
\end{grammarlens}

\subsection*{Your turn}
\begin{itemize}
\raggedright
  \item \emph{Say it:} \emph{disṇe}
  \item \emph{Say it:} \emph{disṇe}, once more
  \item \emph{Say it:} \emph{khiḍkītūn malā ek vicitra prakāś dislā.} — I saw a strange light through the window.
  \item \emph{Recall:} say the Marathi for to flee, then the Marathi for to shut in, then say \emph{disṇe} again
\end{itemize}

\begin{tcolorbox}[breakable,colback=teal!4,colframe=teal!35!black,title={Before you move on}]
What does \mr{दिसणे} mean? (\textquotedblleft{}To be seen; to appear\textquotedblright{}.) Say it once more.
\end{tcolorbox}
\section[\texorpdfstring{\mr{सापडणे} (sāpaḍṇe)}{सापडणे (sāpaḍṇe)}]{\mr{सापडणे} (sāpaḍṇe) — to be found; to turn up}

\label{lesson:MR-C301-sapadne}



\begin{quote}
Before the new one: say the Marathi for to shut in, then the Marathi for to be seen.
\end{quote}

\subsection*{You'll want to know: \mr{सापडणे}}
\textbf{\mr{सापडणे}} — \emph{sāpaḍṇe} — \textquotedblleft{}to be found; to turn up\textquotedblright{}.

\begin{quote}
\mr{बागेत} \mr{मला} \mr{एक} \mr{सोन्याची} \mr{अंगठी} \mr{सापडली}.
\end{quote}

\begin{quote}
\emph{bāget malā ek sonyācī aṅgṭhī sāpaḍlī.}
\end{quote}

\begin{quote}
I found a gold ring in the garden.
\end{quote}

Past (to him …): \mr{त्याला} \mr{सापडलं}, \emph{tyālā sāpaḍlaṁ}.

A dative-subject verb like \mr{दिसणे}: \mr{मला} \mr{किल्ली} \mr{सापडली}, I found the key; \mr{चोर} \mr{सापडला}, the thief was caught.

\begin{cousinweb}
\textbf{\mr{सापडणे}} (\emph{sāpaḍṇe}) is a tadbhava: CDIAL lists Marathi \emph{sāpaḍṇe}, 'to fall into, be got hold of', under Sanskrit \emph{saṁpatati}, 'flies together, reaches, happens' (CDIAL 12930).
\end{cousinweb}

\begin{grammarlens}[title={What you've built}]
One more word for telling what happened.
\end{grammarlens}

\subsection*{Your turn}
\begin{itemize}
\raggedright
  \item \emph{Say it:} \emph{sāpaḍṇe}
  \item \emph{Say it:} \emph{sāpaḍṇe}, once more
  \item \emph{Say it:} \emph{bāget malā ek sonyācī aṅgṭhī sāpaḍlī.} — I found a gold ring in the garden.
  \item \emph{Recall:} say the Marathi for to shut in, then the Marathi for to be seen, then say \emph{sāpaḍṇe} again
\end{itemize}

\begin{tcolorbox}[breakable,colback=teal!4,colframe=teal!35!black,title={Before you move on}]
What does \mr{सापडणे} mean? (\textquotedblleft{}To be found; to turn up\textquotedblright{}.) Say it once more.
\end{tcolorbox}
\section[\texorpdfstring{\mr{आढळणे} (āḍhaḷṇe)}{आढळणे (āḍhaḷṇe)}]{\mr{आढळणे} (āḍhaḷṇe) — to come across}

\label{lesson:MR-C301-adhalne}



\begin{quote}
Before the new one: say the Marathi for to be seen, then the Marathi for to be found.
\end{quote}

\subsection*{You'll want to know: \mr{आढळणे}}
\textbf{\mr{आढळणे}} — \emph{āḍhaḷṇe} — \textquotedblleft{}to come across\textquotedblright{}.

\begin{quote}
\mr{जंगलात} \mr{अभ्यासकांना} \mr{एक} \mr{दुर्मीळ} \mr{पक्षी} \mr{आढळला}.
\end{quote}

\begin{quote}
\emph{jaṅglāt abhyāskānnā ek durmīḷ pakṣī āḍhaḷlā.}
\end{quote}

\begin{quote}
The researchers came across a rare bird in the forest.
\end{quote}

Past (to him …): \mr{त्याला} \mr{आढळलं}, \emph{tyālā āḍhaḷlaṁ}.

A dative subject, more formal than \mr{सापडणे}: \mr{लोकांना} \mr{आढळलं}, people came across it.

\begin{cousinweb}
\textbf{\mr{आढळणे}} (\emph{āḍhaḷṇe}) is inherited: CDIAL lists Marathi \emph{āḍhaḷṇe}, 'to fall in the way of, to meet', under \emph{*āḍḍhalati}, 'falls towards' (CDIAL 1105).
\end{cousinweb}

\begin{grammarlens}[title={What you've built}]
One more word for telling what happened.
\end{grammarlens}

\subsection*{Your turn}
\begin{itemize}
\raggedright
  \item \emph{Say it:} \emph{āḍhaḷṇe}
  \item \emph{Say it:} \emph{āḍhaḷṇe}, once more
  \item \emph{Say it:} \emph{jaṅglāt abhyāskānnā ek durmīḷ pakṣī āḍhaḷlā.} — The researchers came across a rare bird in the forest.
  \item \emph{Recall:} say the Marathi for to be seen, then the Marathi for to be found, then say \emph{āḍhaḷṇe} again
\end{itemize}

\begin{tcolorbox}[breakable,colback=teal!4,colframe=teal!35!black,title={Before you move on}]
What does \mr{आढळणे} mean? (\textquotedblleft{}To come across\textquotedblright{}.) Say it once more.
\end{tcolorbox}
\section[\texorpdfstring{\mr{डोकावणे} (ḍokāvṇe)}{डोकावणे (ḍokāvṇe)}]{\mr{डोकावणे} (ḍokāvṇe) — to peep; to peek}

\label{lesson:MR-C301-dokavne}



\begin{quote}
Before the new one: say the Marathi for to be found, then the Marathi for to come across.
\end{quote}

\subsection*{You'll want to know: \mr{डोकावणे}}
\textbf{\mr{डोकावणे}} — \emph{ḍokāvṇe} — \textquotedblleft{}to peep; to peek\textquotedblright{}.

\begin{quote}
\mr{मुलगा} \mr{दारातून} \mr{आत} \mr{डोकावला}.
\end{quote}

\begin{quote}
\emph{mulgā dārātūn āt ḍokāvlā.}
\end{quote}

\begin{quote}
The boy peeped in through the doorway.
\end{quote}

Past (he …): \mr{तो} \mr{डोकावला}, \emph{to ḍokāvlā}.

Built on \mr{डोके}, the head: to put one's head in for a look.

\begin{cousinweb}
\textbf{\mr{डोकावणे}} (\emph{ḍokāvṇe}) is built on \emph{ḍoke}, 'the head', which goes back to \emph{*ḍōka}, 'head' (CDIAL 5566): to push the head in for a look.
\end{cousinweb}

\begin{grammarlens}[title={What you've built}]
One more word for telling what happened.
\end{grammarlens}

\subsection*{Your turn}
\begin{itemize}
\raggedright
  \item \emph{Say it:} \emph{ḍokāvṇe}
  \item \emph{Say it:} \emph{ḍokāvṇe}, once more
  \item \emph{Say it:} \emph{mulgā dārātūn āt ḍokāvlā.} — The boy peeped in through the doorway.
  \item \emph{Recall:} say the Marathi for to be found, then the Marathi for to come across, then say \emph{ḍokāvṇe} again
\end{itemize}

\begin{tcolorbox}[breakable,colback=teal!4,colframe=teal!35!black,title={Before you move on}]
What does \mr{डोकावणे} mean? (\textquotedblleft{}To peep; to peek\textquotedblright{}.) Say it once more.
\end{tcolorbox}
\section[\texorpdfstring{\mr{निरखणे} (nirakhṇe)}{निरखणे (nirakhṇe)}]{\mr{निरखणे} (nirakhṇe) — to observe closely}

\label{lesson:MR-C301-nirakhne}



\begin{quote}
Before the new one: say the Marathi for to come across, then the Marathi for to peep.
\end{quote}

\subsection*{You'll want to know: \mr{निरखणे}}
\textbf{\mr{निरखणे}} — \emph{nirakhṇe} — \textquotedblleft{}to observe closely\textquotedblright{}.

\begin{quote}
\mr{आजोबांनी} \mr{नाणं} \mr{उलटं} \mr{करून} \mr{नीट} \mr{निरखलं}.
\end{quote}

\begin{quote}
\emph{ājobānnī nāṇaṁ ulṭaṁ karūn nīṭ nirakhlaṁ.}
\end{quote}

\begin{quote}
Grandfather turned the coin over and studied it closely.
\end{quote}

Past, with -\mr{ने} (he …): \mr{त्याने} \mr{निरखलं}, \emph{tyāne nirakhlaṁ}.

Transitive: the thing studied is the object. \mr{निरीक्षण} \mr{करणे} is its formal cousin.

\begin{cousinweb}
\textbf{\mr{निरखणे}} (\emph{nirakhṇe}) is a tadbhava: CDIAL lists Marathi \emph{nirakhṇe}, of close looking, under Sanskrit \emph{nirīkṣate}, 'looks at, regards' (CDIAL 7280).
\end{cousinweb}

\begin{grammarlens}[title={What you've built}]
One more word for telling what happened.
\end{grammarlens}

\subsection*{Your turn}
\begin{itemize}
\raggedright
  \item \emph{Say it:} \emph{nirakhṇe}
  \item \emph{Say it:} \emph{nirakhṇe}, once more
  \item \emph{Say it:} \emph{ājobānnī nāṇaṁ ulṭaṁ karūn nīṭ nirakhlaṁ.} — Grandfather turned the coin over and studied it closely.
  \item \emph{Recall:} say the Marathi for to come across, then the Marathi for to peep, then say \emph{nirakhṇe} again
\end{itemize}

\begin{tcolorbox}[breakable,colback=teal!4,colframe=teal!35!black,title={Before you move on}]
What does \mr{निरखणे} mean? (\textquotedblleft{}To observe closely\textquotedblright{}.) Say it once more.
\end{tcolorbox}
